\documentclass{expreport}
% =============================================
% Part 0 Edit the info
% =============================================
\usepackage{multirow}
\usepackage{array}
\usepackage{booktabs}
\usepackage{caption}
\usepackage{subfigure}
\title{东南大学电工电子实验中心}
\course{coursename}
\college{collegename}
\date{2022年x月x日}
\lab{labname}
\expname{expname}
\name{name}
\major{major}
\stuid{id}
\group{group}
\partner{partner}
\ttitle{第\ x\ 次实验}

\begin{document}
% =============================================
% Part 1 Header
% =============================================
\makecover
% \makeheader

% =============================================
% Part 2 Main document
% =============================================

\section{实验目的}

    \subsection{副标题} 
   
 

\section{实验原理}

\subsection{副标题}

\section{实验内容}


\subsection{副标题}


\begin{table}[H]
    \begin{center}
     \setlength{\tabcolsep}{0.17cm}{
        \begin{tabular}{|l|l|l|l|l|l|l|l|l|l|l|l|l|l|}
            \hline
            探头衰减 & \multicolumn{3}{|c|}{幅度} & \multicolumn{3}{|c|}{高电平电压} & \multicolumn{3}{|c|}{低电平电压} & \multicolumn{3}{|c|}{周期} & 频率  \\
            \hline
            \multirow{2}*{$\times$1} & 档位& 格数 &计数值 & 档位& 格数 &计数值& 档位& 格数 &计数值& 档位& 格数 &计数值 &频率\\
         \cline{2-14}
         &  500 & 6 & 3.0 & 500 & 6 & 3.0 &500 & 0&0 &200 & 5 &$10^3$ & $10^3$ \\
          \hline
       
        \end{tabular}}
    \end{center}
    \caption{示波器补偿（校准）信号测量}
    \label{fig:dist}
\end{table}

波形记录：
\begin{figure}[H]
              \centering
              \includegraphics[width=0.6\linewidth]{figures/01.jpg}
              \caption{补偿信号测量}
              \label{fig:dist}
          \end{figure}
   


\begin{figure}[H]
  \centering
  \subfigure[峰峰值]{
    \label{fig:subfig:onefunction} 
    \includegraphics[width=0.4\linewidth]{figures/01.jpg}}
  \hspace{0.5in} % 两图片之间的距离
  \subfigure[周期]{
    \label{fig:subfig:twofunction} 
    \includegraphics[width=0.4\linewidth]{figures/01.jpg}}
  \label{fig:dist} 
\end{figure}

\section{实验使用仪器设备}


\section{实验总结}
\subsection{误差分析}

\subsection{收获领悟}

\section{参考资料}

\end{document}
